\documentclass[letterpaper,11pt,twoside]{article}
\usepackage[utf8]{inputenc}
\usepackage{amsmath,amsfonts,amssymb,amsthm,latexsym}
\usepackage[spanish,es-noshorthands]{babel}
\usepackage[T1]{fontenc}
\usepackage{lmodern}
\usepackage{graphicx,hyperref}
\usepackage{tikz,pgf}
\usepackage{multicol}
\usepackage{fancyhdr}
\usepackage[height=9.5in,width=7in]{geometry}
\usepackage{fancyhdr}
\pagestyle{fancy}
\fancyhead[LE]{\includegraphics[height=12pt]{Images/logo-colegio.png} Probabilidad $11^{\circ}$}
\fancyhead[RE]{}
\fancyhead[RO]{\textit{Germ\'an Avenda\~no Ram\'irez, Lic. U.D., M.Sc. U.N.}}
\fancyhead[LO]{}

\author{Germ\'an Avenda\~no Ram\'irez, Lic. U.D., M.Sc. U.N.}
\title{\begin{minipage}{.2\textwidth}
\includegraphics[height=1.75cm]{Images/logo-colegio.png}\end{minipage}
\begin{minipage}{.55\textwidth}
\begin{center}
Taller, Probabilidad eventos empíricos\\
Probabilidad $11^{\circ}$
\end{center}
\end{minipage}\hfill
\begin{minipage}{.2\textwidth}
\includegraphics[height=1.75cm]{Images/logo-sed.png} 
\end{minipage}}
\date{}
\thispagestyle{plain}
\begin{document}
\maketitle
Nombre: \hrulefill Curso: \underline{\hspace*{44pt}} Fecha: \underline{\hspace*{2.5cm}}
\begin{multicols}{2}
 \begin{enumerate}
  \item Si usted lanza un dado 40 veces y 9 de los tiros resultan en un ``5'', ¿qué probabilidad empírica se observó para el evento ``5''?
  \item Explique por qué una probabilidad empírica, una proporción observada, y una frecuencia relativa son en realidad tres nombres diferentes para lo mismo.
  \item Millones de personas viajan en ferrocarril todos los años. La Asociación Nacional de Pasajeros de Ferrocarril proporciona las siguientes cantidades de viajes en 2004.
  \begin{center}
\begin{tabular}{ll}
Sistema ferroviario & Viajeros (millones)\\ \hline
Sistema Amtrak & 25.0\\
Corredor Noroeste & 14.2\\
Suburbano + 0este & 10.8 \hline
  \end{tabular}
  \end{center}
\begin{enumerate}
 \item ¿Qué porcentaje de pasajeros de ferrocarril usaron
el sistema Amtrak en 2004?
\item Si uno de estos pasajeros ha de ser entrevistado, ¿cuál es la probabilidad de que él haya viajado en el sistema Amtrak en 2004 si es seleccionado al azar?
\item Explique la diferencia y la relación entre preguntas y respuestaas  a las partes \textit{a} y \textit{b}.
\end{enumerate}

 \end{enumerate}

\end{multicols}


\end{document}
